\section{Discusión}
\textbf{Acuerdo alcanzado.} El clasificador por reglas llegó a $\kappa = 0{,}343$ sobre los pacientes de prueba, lejos del acuerdo entre los técnicos del HMC (0,74) y de la red de referencia (0,79). Es un acuerdo bajo pero claramente superior al azar, y conviene leerlo junto con lo que cada número representa: la red aprende de miles de épocas etiquetadas qué combinación de rasgos separa los estadios, mientras que acá cada estadio se reconoce con dos o tres descriptores y umbrales fijos, que se pueden leer y justificar uno por uno. El $\kappa$ además varió mucho entre pacientes (de 0,16 a 0,55, con mediana 0,28), lo que indica que parte del error depende de la calidad y de las particularidades de cada registro.

\textbf{Dónde se equivoca.} Los errores se concentran en las fronteras que también son difíciles para los técnicos. N1 fue el estadio peor reconocido (F1 de 0,18): se asigna por descarte, de modo que recibe las épocas de cualquier estadio que no cumplen ninguna regla, y solo el 16\,\% de las épocas asignadas así eran realmente N1. La frontera entre N2 y N3 es gradual, porque N3 se define por la proporción de la época ocupada por ondas lentas; esto se vio al elegir el umbral de ondas, donde 2 y 3 dieron casi el mismo acuerdo, uno a costa de N2 y el otro a costa de N3. Por último, las épocas N1 con movimientos oculares y tono muscular conservado se parecen a la vigilia con ojos abiertos, que la regla de W reconoce por los parpadeos.

\textbf{Suavizado.} Las reglas de continuidad mejoraron poco el acuerdo (de 0,330 a 0,343), pero con un balance favorable: por cada error que introdujeron corrigieron 1,4. El filtro de moda cambió casi el doble de épocas con un balance peor, porque también borra tramos cortos reales, como los despertares breves. Que el suavizado aporte poco indica que la mayoría de los errores no son épocas aisladas, sino tramos enteros asignados al estadio equivocado.

\textbf{Reglas frente a LDA + HMM.} Al reemplazar el LDA y el HMM por reglas fijas se resignó parte del acuerdo esperable: un clasificador supervisado pondera todos los descriptores a la vez, mientras que cada regla usa solo dos o tres. A cambio, cada época queda explicada por una regla concreta (Fig.~\ref{fig:uso_reglas}), con umbrales que tienen interpretación fisiológica y que no dependen de un entrenamiento. Esa trazabilidad permitió diagnosticar cada error: la regla de N3 acierta menos (44\,\%) porque cuenta como ondas lentas las deflexiones de los movimientos, y la de N1 casi nunca acierta porque recibe lo que las demás descartan. Una primera versión, que expresaba todos los descriptores como percentiles de la noche, mostró además que normalizar por paciente sin criterio es contraproducente: fijaba la proporción de cada estadio en cada noche y llevaba a N3 épocas de vigilia de pacientes con poco sueño profundo.

\textbf{Preprocesamiento.} Cada parámetro del preprocesamiento salió de una medición sobre los pacientes de entrenamiento, y dos de esas mediciones cambiaron lo que proponía el anteproyecto. El pasa-altos tuvo que bajar a \SI{0,2}{\hertz} para no perder ondas lentas. Y la corrección ocular aplicada a toda la noche borraba ondas lentas reales en N3, porque E1$-$E2 también registra actividad cerebral, lo que obligó a agregar la compuerta. La remoción cardíaca, que la cátedra puso en duda, resultó relevante por un motivo que no era evidente. El artefacto tiene la misma amplitud toda la noche, pero pesa más en REM, donde el EEG es más chico, y sin quitarlo REM se parece más a W y a N2 en las bandas alfa y sigma.

La marca de no estacionariedad mostró una fuente concreta de error del clasificador. Un tercio de las épocas marcadas de W, N1, N2 y REM se clasificó como N3, porque los movimientos y los despertares producen deflexiones grandes y lentas que la regla de N3 cuenta como ondas lentas. Una mejora posible es no asignar N3 a las épocas marcadas, o contar las ondas solo en los tramos estables de la época.

El preprocesamiento tiene cuatro limitaciones. El equipo aplicó un pasa-bajos de \SI{35}{\hertz} en 27 de los 56 pacientes de entrenamiento y de \SI{128}{\hertz} en el resto, lo que puede correr la pendiente espectral de un grupo respecto del otro. SN020 tiene más de un tercio de sus épocas saturadas, y la HRV de SN001 es poco confiable por su ritmo irregular. Por último, en W se siguen detectando ondas lentas que probablemente sean movimientos del cuerpo que el EOG no registra.

\textbf{Descriptores y técnicas no vistas en clase.} H1 se confirmó, porque el multitaper varió menos que Welch entre épocas vecinas de un mismo estadio en los 14 pacientes de prueba (Tabla~\ref{tab:hipotesis}). La mejora fue del 18\,\% al comparar frecuencia por frecuencia, pero solo del 4\,\% al comparar potencias por banda, ya que integrar una banda promedia muchas frecuencias y Welch se acerca así al multitaper. Como los descriptores son potencias por banda, la ventaja del multitaper existe pero es chica para este uso.

H2 se confirmó con un margen estrecho. Con ondas de forma conocida, Woody tuvo menos error que el promedio ordinario en 11 de los 14 pacientes ($p = 0{,}049$). En las ondas reales el promedio ordinario dio un valle más profundo, pero eso no indica una forma más fiel, sino que cada recorte está centrado en su propio mínimo y el ruido se suma justo ahí. Que el algoritmo no haya convergido en más de la mitad de los pacientes sugiere agregar una tolerancia, de modo que termine cuando cambian pocas ondas y no cuando no cambia ninguna. La segunda parte de H2, que Woody ayude a separar N2 de N3, no se sostiene. La similitud con la plantilla es casi igual en todos los estadios, y lo que separa N3 es la cantidad de ondas y no su forma.

H3 no se confirmó. El desempate con CSP mejoró levemente el F1 de N1 (mediana de 0,16 frente a 0,14; fue peor con CSP solo en 5 de los 14 pacientes), sin que la diferencia fuera significativa ($p = 0{,}060$), y empeoró el de REM (0,46 frente a 0,57, peor en 10 de los 14; $p = 0{,}009$), por lo que el clasificador final no lo usa. En un paciente aislado CSP separa N1 de REM, pero los filtros se ajustaron con las covarianzas de los 56 pacientes de entrenamiento juntos, y la combinación de canales que separa los dos estadios cambia de un paciente a otro. Además, la clase N1 que corrige el desempate es la que asignan las reglas por descarte, que mezcla épocas de todos los estadios.

El hallazgo más útil de los descriptores fue la vigilia activa. Las épocas de W con movimiento tienen EMG alto, poca $\alpha$, y también husos y ondas lentas detectados. Sin $\alpha$ no cumplían la regla de W, y terminaban en N3 por las ondas, en N2 por los husos o en N1 por descarte. Por eso la regla final de W acepta, en lugar de la $\alpha$, parpadeos (EOG alto), que recuperan la vigilia con ojos abiertos. La vigilia con movimiento pero sin parpadeos sigue sin reconocerse; una regla que acepte un EMG muy alto aunque falten la $\alpha$ y los parpadeos atacaría ese error restante.

H4 se confirmó solo en parte. La remoción eliminó el artefacto cardíaco en los 14 pacientes de prueba, con un residuo mediano de 1,2~\si{\micro\volt} frente a 19,5~\si{\micro\volt} antes de remover ($p = 0{,}0001$), pero el acuerdo con los técnicos no cambió al quitar los artefactos y la HRV ($\kappa$ mediano de 0,28 con estas etapas y 0,29 sin ellas; fue menor con las etapas en 6 de los 14 pacientes, $p = 0{,}54$). El artefacto existe y la remoción funciona, como mostró también el análisis del preprocesamiento, pero las reglas finales casi no lo perciben: usan potencias relativas sobre épocas de 30~s, en las que un pulso de pocos microvoltios pesa poco, y el EMG y el EOG, que se registran en canales aparte. El cociente LF/HF, por su parte, separó poco REM de N1 y N2. Para este clasificador, entonces, la remoción del artefacto cardíaco no era necesaria, como sugirió la devolución; sí lo sería para descriptores más sensibles a la forma de onda.

\textbf{Limitaciones.} Los umbrales se eligieron con una sola partición en 56 pacientes de entrenamiento y 14 de prueba; con tan pocos pacientes de prueba, el $\kappa$ global depende bastante de cuáles tocaron en cada grupo. Durante el desarrollo, el acuerdo sobre los pacientes de prueba se consultó para descartar la primera versión basada en percentiles y en un ajuste intermedio del umbral de ondas lentas; ese ajuste se revirtió y los umbrales finales se eligieron solo con los pacientes de entrenamiento, pero el resultado de prueba no es completamente independiente del diseño. Por último, en esta base casi no hay épocas fuera del período entre luces apagadas y encendidas, por lo que esa decisión no pudo evaluarse.
